\subsection{40. DRAS: Deep Reinforcement Learning for Cluster Scheduling in HPC}
\textbf{OSTI:} \texttt{1984484}\quad\textbf{Authors:} Fan et al.\ (2021)\quad\textbf{Domain:} Systems / Reinforcement Learning\\
\scorebadge{Coverage}{8}\quad\scorebadge{Agreement}{8}\quad\textbf{Total:} 16/20\par\smallskip
\noindent\textbf{Coverage rationale.} Built an event-driven HPC cluster simulator with DQN and PPO agents trained on the HPC2N workload trace (240 nodes, 5,000 jobs). Compared against FCFS, EASY-backfill, and SJF baselines. PPO achieves 24\% lower average slowdown than SJF and 81\% lower than FCFS, supporting the paper's core DRL-over-traditional-scheduling claim.\par
\noindent\textbf{Agreement rationale.} Paper claims $\sim$20--30\% DRL improvement over best traditional scheduler; we observe 24\% PPO vs.\ SJF. Ranking: PPO $>$ DQN $>$ SJF $>$ EASY-BF $>$ FCFS matches paper Table~3 order. Makespan improvement within 0.3\% of SJF (paper also shows near-parity on makespan).\par\smallskip
\noindent\textbf{What reproduced:}
\begin{itemize}[leftmargin=1.4em,itemsep=0.2em,topsep=0.2em]
\item Event-driven simulator on HPC2N trace (240 nodes, 5K jobs)
\item FCFS, EASY-BF, SJF baseline schedulers
\item DQN and PPO agents with experience replay
\item PPO avg slowdown 125.8 vs.\ SJF 166.3 (24\% improvement)
\item Ranking: PPO $>$ DQN $>$ SJF $>$ EASY-BF $>$ FCFS
\end{itemize}
\noindent\textbf{What missing:}
\begin{itemize}[leftmargin=1.4em,itemsep=0.2em,topsep=0.2em]
\item Paper's hierarchical two-network DRAS (we used single-network DQN/PPO)
\item SDSC Bluehorizon and ANL traces (used only HPC2N)
\item Long-horizon job completion time vs.\ paper's reported values
\item Multi-seed statistical error bars (single seed)
\end{itemize}
\noindent\textbf{Follow-on research questions:}
\begin{enumerate}[leftmargin=1.6em,itemsep=0.2em,topsep=0.2em,label=Q\arabic*.]
\item Does the hierarchical two-network DRAS architecture produce measurably better slowdown than single-network PPO, and what is the added training cost?
\item Evaluate DRL scheduler generalization: train on HPC2N, test on SDSC BH and ANL traces without fine-tuning --- what is the zero-shot performance gap?
\item Add heterogeneous node types (CPU-only vs.\ GPU-enabled) and measure how PPO's state representation handles resource fragmentation.
\item Apply multi-objective RL (slowdown + energy cost) and trace the Pareto frontier --- does the energy term conflict with makespan minimization under backfill?
\item Compare DRAS against a modern graph-based policy (GNN over job dependency graphs) on the HPC2N trace --- does topological job-structure awareness improve throughput beyond the DRAS improvement?
\end{enumerate}

